\section{\qI}

\subsection{Introduction}
This subquestion investigates the mathematical structure of the I - V - IV - iv chord progression, by calculating the normalized harmonic complexity of it and other chord progressions, and comparing the results.

\subsection{Results}
The following is the standard output of my Python script, which performs the calculations for the research paper [maybe no stdout but a table]:
\begin{lstlisting}
    Ratio major (4, 5, 6)
    Ratio minor (10, 12, 15)
    Ratio octave (1, 2, 4)
    Ratio adjacent notes (24, 25, 27)
    Normalized harmonic complexity of octave =  0.0
    Normalized harmonic complexity of I_V_IV_IV =  35.57
    Normalized harmonic complexity of ! I_V_IV_iv =  44.46
    Normalized harmonic complexity of I_V_iv_iv =  53.35
    Normalized harmonic complexity of I_v_iv_iv =  62.25
    Normalized harmonic complexity of i_v_iv_iv =  71.14
    Normalized harmonic complexity of adj =  100.0
    Normalized harmonic complexity of (popular prog) I_V_vi_IV =  44.46
    Normalized harmonic complexity of (popular prog) I_V_vi_iii_IV =  49.8
    Normalized harmonic complexity of (popular prog) vi_V_IV_V =  44.46
    Arc of harmonic complexity of octave =  0.0
    Arc of harmonic complexity of I_V_IV_IV =  35.57 - 35.57 - 35.57 - 35.57
    Arc of harmonic complexity of ! I_V_IV_iv =  35.57 - 35.57 - 35.57 - 71.14
    Arc of harmonic complexity of I_V_iv_iv =  35.57 - 35.57 - 71.14 - 71.14
    Arc of harmonic complexity of I_v_iv_iv =  35.57 - 71.14 - 71.14 - 71.14
    Arc of harmonic complexity of i_v_iv_iv =  71.14 - 71.14 - 71.14 - 71.14
    Arc of harmonic complexity of adj =  100.0
    Arc of harmonic complexity of (popular prog) I_V_vi_IV =  35.57 - 35.57 - 71.14 - 35.57
    Arc of harmonic complexity of (popular prog) I_V_vi_iii_IV =  35.57 - 35.57 - 71.14 - 71.14 - 35.57
    Arc of harmonic complexity of (popular prog) vi_V_IV_V =  71.14 - 35.57 - 35.57 - 35.57
\end{lstlisting}

% The following are the matplotlib oiutputs:
% standard dev of harmonic complexity of popular progressions
% spectrum of harmonic complexity of chord progressions

\subsection{Analysis}
The calculations yield that the I - V - IV - iv chord progression has a normalized harmonic complexity of ~44.46\%, which is similar to . A chord progression with solely major chords — like I - V - IV - IV — is closer to this 45\% spot of normalized harmonic complexity than a chord progression with solely minor chords — like i - v - iv - iv. [No reliable external data could be found on the (un)popularity of all chord progressions in the list, only that the progressions labeled "(popular prog)" are commonly used chord progressions in music; therefore, (un)popularity of each chord progressions will not be compared.]

External research suggests that perceived enjoyment of music in terms of harmonic complexity (y = perceived enjoyment and x = amount of harmonic complexity) has an inverted U shape, also known as the Wundt Curve \ref{Wundt Curve resource}. [The peak of this curve would supposedly be reached at ~50\% of harmonic complexity, where 0\% is very low complexity, and 100\% is extreme complexity.](ja? Is dat zo? Ik herinner niet van waar ik dit heb. Change it to this: we cannot conclude exatcly where the chord progression is. En dan koppelen aan subquestion 2)

\subsection{Conclusion}
Given that the peak of perceived enjoyment should be at ~50\% harmonic complexity, and that the I - V - IV - iv chord progression lies at ~45\% harmonic complexity, we can infer that the I - V - IV - iv chord progression sits in a "sweet spot" of normalized harmonic complexity. [This is similar to other popular chord progressions which have the same amount of major/minor chords, but distinct from chord progressions ranging in the amount of major/minor chords.]